\documentclass[12pt]{article}
\usepackage[utf8]{inputenc}
\usepackage[spanish]{babel}
\usepackage[a4paper, margin=2cm]{geometry}
\usepackage{tikz}
\usetikzlibrary{babel}
\usepackage[most]{tcolorbox}
\usepackage{amsmath}

% Usar fuente sin serifa para un aspecto más moderno (tipo infografía)
\renewcommand{\familydefault}{\sfdefault}

% Definición de colores con temática médica/científica
\definecolor{medblue}{RGB}{43, 108, 176}
\definecolor{medteal}{RGB}{49, 151, 149}
\definecolor{medlight}{RGB}{235, 248, 255}
\definecolor{meddark}{RGB}{26, 32, 44}
\definecolor{tracinitial}{RGB}{113, 128, 150} % Gris azulado para vector inicial
\definecolor{tracfinal}{RGB}{49, 130, 206}    % Azul brillante para vector final
\definecolor{tracdelta}{RGB}{229, 62, 62}     % Rojo para la diferencia/cambio

\begin{document}
\pagestyle{empty}

\begin{tcolorbox}[
    enhanced, % Habilita características avanzadas como las sombras
    colback=medlight!30, 
    colframe=medblue, 
    title={\Large \textbf{Resta de Vectores: Tracción Ortopédica}}, 
    halign title=center, 
    fonttitle=\bfseries, 
    arc=4mm, 
    boxrule=1.5pt,
    shadow={2mm}{-2mm}{0mm}{black!30!white}
]

    \vspace{0.2cm}
    \begin{tcolorbox}[colback=white, colframe=medteal, arc=2mm, boxrule=1.2pt, title=\textbf{Caso Clínico / Problema}]
        Un sistema ortopédico ejercía una tracción inicial $\vec{T}_1 = (50\hat{i} + 20\hat{j})$ N. Después de reajustar los pesos, la tracción es $\vec{T}_2 = (30\hat{i} + 40\hat{j})$ N. Calcule el vector de cambio de fuerza $\Delta\vec{T} = \vec{T}_2 - \vec{T}_1$.
    \end{tcolorbox}

    \vspace{0.4cm}
    
    \begin{center}
    \begin{tikzpicture}[scale=0.12]
        % Cuadrícula de fondo (escala de 10 en 10)
        \draw[step=10,gray!40,very thin] (-35,-5) grid (55,55);
        
        % Ejes cartesianos
        \draw[thick,->,meddark] (-35,0) -- (55,0) node[anchor=north west] {$x$ (N)};
        \draw[thick,->,meddark] (0,-5) -- (0,55) node[anchor=south east] {$y$ (N)};
        
        % Marcas numéricas (eje x e y)
        \foreach \x in {-30,-20,-10,10,20,30,40,50} \draw (\x, 1.5) -- (\x, -1.5) node[anchor=north, font=\scriptsize, yshift=-2pt] {$\x$};
        \foreach \y in {10,20,30,40,50} \draw (1.5, \y) -- (-1.5, \y) node[anchor=east, font=\scriptsize, xshift=-2pt] {$\y$};
        
        % Coordenadas de los vectores
        \coordinate (O) at (0,0);
        \coordinate (T1) at (50,20);
        \coordinate (T2) at (30,40);
        \coordinate (DT) at (-20,20);
        
        % Vector Inicial T1
        \draw[->, line width=2pt, tracinitial] (O) -- (T1) node[midway, below right, font=\bfseries] {$\vec{T}_1$};
        
        % Vector Final T2
        \draw[->, line width=2pt, tracfinal] (O) -- (T2) node[pos=0.6, left, font=\bfseries] {$\vec{T}_2$};
        
        % Vector Cambio (Resta Geométrica: de T1 a T2)
        \draw[->, line width=2.5pt, tracdelta] (T1) -- (T2) node[midway, above right, font=\bfseries, text=tracdelta] {$\Delta\vec{T}$};
        
        % Proyección del Cambio desde el origen
        \draw[->, dashed, line width=2pt, tracdelta!60] (O) -- (DT) node[midway, above left, font=\bfseries, text=tracdelta] {$\Delta\vec{T}$};
        
        % Puntos destacados
        \filldraw[meddark] (O) circle (4pt) node[anchor=north east, font=\bfseries, fill=white, inner sep=1pt] {Origen};
        \filldraw[tracinitial] (T1) circle (4pt) node[anchor=north west, font=\bfseries, text=tracinitial] {$(50,20)$};
        \filldraw[tracfinal] (T2) circle (4pt) node[anchor=south, font=\bfseries, text=tracfinal, yshift=2pt] {$(30,40)$};
        
    \end{tikzpicture}
    \end{center}

    \vspace{0.4cm}
    \begin{tcolorbox}[colback=medteal!10, colframe=medteal, arc=2mm, boxrule=1.2pt, title=\textbf{Desarrollo Matemático y Solución}]
        Para encontrar el vector de cambio $\Delta\vec{T}$, restamos las componentes del vector inicial $\vec{T}_1$ al vector final $\vec{T}_2$:
        \begin{align*}
            \Delta\vec{T} &= \vec{T}_2 - \vec{T}_1 \\[1ex]
            \Delta\vec{T} &= (30\hat{i} + 40\hat{j}) - (50\hat{i} + 20\hat{j}) = (30 - 50)\hat{i} + (40 - 20)\hat{j} \\[1ex]
            \Delta\vec{T} &= \mathbf{-20\hat{i} + 20\hat{j} \text{ N}}
        \end{align*}
        \textbf{Resultado:} El vector de cambio de fuerza reajustado es \textbf{$-20\hat{i} + 20\hat{j}$ N}.
    \end{tcolorbox}
    
\end{tcolorbox}

\end{document}